\chapter{External Device Drivers}
\label{app:external}

This appendix is the companion to appendix~\ref{app:internal}: where that one
covers the SoC's on-chip peripherals, this one covers the drivers for
\emph{external} parts -- the I/O expanders, sensors, clocks, codec, display,
GPS module, SD card, LED string and Ethernet controller that hang off the
ESP32-S3's buses. Each driver layers on one (or two) of the on-chip drivers: the
I2C devices ride \texttt{ESP32S3.I2C}, the display rides \texttt{ESP32S3.SPI}, the
codec uses \texttt{ESP32S3.I2C} for control and \texttt{ESP32S3.I2S} for audio, the
GPS service reads \texttt{ESP32S3.UART}, and the LED string is clocked out over
\texttt{ESP32S3.RMT}. The W5500 Ethernet controller is the largest of these -- a
whole networking stack -- so it has its own chapter
(\ref{ch:networking}); this appendix gives only its bus and a pointer.

The same ownership discipline recurs. Most of these drivers expose a
\texttt{Device} configured once by \texttt{Setup} and, for the multi-step ones, a
\texttt{Session} (or \texttt{Output}/\texttt{Strip}) handle taken with
\texttt{Acquire} that holds the part exclusively across a sequence of operations
and auto-releases on scope exit, while the underlying bus is locked only per
transaction -- so other devices on the same bus stay reachable between your
calls. Each entry gives the device's bus, its I2C address where relevant, a digest
of the public API, and a worked example from the \texttt{examples/} tree. The
controlled handles mean these drivers require the \texttt{embedded} or
\texttt{full} profile, as noted per section.

\section{CH422G -- I2C I/O Expander}\index{CH422G driver}
\label{app:ext:ch422g}

The WCH CH422G is an I2C I/O expander offering eight bidirectional pins
(\texttt{IO0..IO7}) plus four output-only pins (\texttt{OC0..OC3}). Unlike a
conventional register-pointer expander, it is driven entirely through \emph{fixed,
function-specific} I2C command addresses -- \texttt{0x24} (write system config),
\texttt{0x23} (write OC outputs), \texttt{0x38} (write IO outputs) and \texttt{0x26}
(read IO inputs) -- so there are no address straps and \texttt{Setup} takes no
address: one CH422G per bus.

\texttt{ESP32S3.CH422G} sits on \texttt{ESP32S3.I2C} and follows the two-level
locking shape of the other expander/RTC drivers: \texttt{Acquire} a
\texttt{Session} to hold the chip across many operations, while the I2C host is
locked only per transaction. Direction is \emph{global} -- one
\texttt{IO\_Direction} bit makes all of \texttt{IO0..IO7} inputs or outputs;
\texttt{OC0..OC3} are output-only with a global \texttt{Push\_Pull}/%
\texttt{Open\_Drain} drive. Because the config/OC/IO-output registers cannot be
read back, the driver keeps a shadow initialised to the datasheet power-on
defaults (IO = inputs, OC = high, push-pull). The chip has no interrupt output,
hence no \texttt{.Interrupts} child. The controlled \texttt{Session} restricts this
driver to the embedded and full profiles.

\begin{lstlisting}
type IO_Value is mod 2 ** 8;          --  bit i = IOi
type OC_Value is mod 2 ** 4;          --  low nibble = OC0..OC3
type IO_Pin is range 0 .. 7;  type OC_Pin is range 0 .. 3;
type Pin_State    is (Low, High);
type IO_Direction is (Inputs, Outputs);        --  all of IO0..IO7
type OC_Drive     is (Push_Pull, Open_Drain);  --  all of OC0..OC3
type Status  is (OK, Bus_Error);
type Device  is limited private;
type Session is limited private;

procedure Setup (Dev : out Device; Sda, Scl : ESP32S3.GPIO.Pin_Id;
                 Host : ESP32S3.I2C.I2C_Host := ESP32S3.I2C.I2C0;
                 Clock_Hz : Positive := 400_000);
procedure Acquire (S : in out Session; Dev : Device);
procedure Release (S : in out Session);
function  Present (S : Session) return Boolean;   --  ACK probe of 0x24
procedure Configure (S : Session; IO_Dir : IO_Direction := Inputs;
                     OC_Mode : OC_Drive := Push_Pull; Result : out Status);
procedure Write_IO  (S : Session; Value : IO_Value; Result : out Status);
procedure Read_IO   (S : Session; Value : out IO_Value; Result : out Status);
procedure Write_OC  (S : Session; Value : OC_Value; Result : out Status);
\end{lstlisting}

The worked example sets the chip up on I2C0 (SDA=IO8, SCL=IO9), probes the config
command address, then reads the input pins once a second. It is deliberately
read-only and never drives a pin.

\begin{lstlisting}
package CH renames ESP32S3.CH422G;
use type CH.Status;
Dev : CH.Device;  S : CH.Session;  V : CH.IO_Value;  St : CH.Status;
begin
   CH.Setup (Dev, Sda => 8, Scl => 9);   --  I2C0, 400 kHz, no address
   CH.Acquire (S, Dev);                  --  held for the whole run

   Put_Line (if CH.Present (S) then "ACK (present)" else "no ACK");

   loop
      delay until Clock + Seconds (1);
      CH.Read_IO (S, V, St);             --  RD-IO (0x26)
      if St = CH.OK then
         Put ("IO inputs = 0x");  Put_Hex (Unsigned_32 (V), 2);  New_Line;
      end if;
   end loop;
\end{lstlisting}

\noindent\textbf{Example:} \texttt{examples/esp32s3\_ch422g}.

\section{TCA9555 -- I2C 16-bit I/O Expander}\index{TCA9555 driver}
\label{app:ext:tca9555}

The Texas Instruments TCA9555 is a 16-bit I2C GPIO expander: two 8-bit ports
(\texttt{P0} = pins 0..7, \texttt{P1} = pins 8..15), each pin independently an
input or output, with an optional input-polarity-inversion register and an
active-low open-drain \texttt{INT} output. The 7-bit address is
\texttt{Base\_Address} (\texttt{0x20}) plus the \texttt{A2/A1/A0} strap value
(0..7), so up to eight parts share one bus.

\texttt{ESP32S3.TCA9555} layers on \texttt{ESP32S3.I2C} and is the canonical
two-level-locking driver: a \texttt{Session} is an exclusive RAII hold on one
expander (so a per-pin read-modify-write is safe), while the I2C host is locked
only inside each transaction -- leaving the bus free for other devices between
your operations. The per-device locks are a fixed library-level array keyed by
(host, strap). The optional \texttt{INT} pin is wired through \texttt{Setup}'s
\texttt{Int\_Pin}, and the \texttt{ESP32S3.TCA9555.Interrupts} child arms it as a
pulled-up falling-edge input. The controlled \texttt{Session} restricts this
driver to the embedded and full profiles.

\begin{lstlisting}
Base_Address : constant ESP32S3.I2C.Slave_Address := 16#20#;
subtype Hardware_Address is Natural range 0 .. 7;
type Pin_Number is range 0 .. 15;     --  0..7 = Port 0, 8..15 = Port 1
type Port_Value is mod 2 ** 16;       --  bit i = Pin_Number i
type Direction is (Output, Input);    type Pin_State is (Low, High);
type Status is (OK, Bus_Error);
type Device is limited private;  type Session is limited private;

procedure Setup (Dev : out Device; Addr : Hardware_Address;
                 Sda, Scl : ESP32S3.GPIO.Pin_Id;
                 Int_Pin  : ESP32S3.GPIO.Optional_Pin := ESP32S3.GPIO.No_Pin;
                 Host : ESP32S3.I2C.I2C_Host := ESP32S3.I2C.I2C0;
                 Clock_Hz : Positive := 400_000);
procedure Acquire (S : in out Session; Dev : Device);
procedure Release (S : in out Session);
procedure Set_Directions (S : Session; Inputs : Port_Value; Result : out Status);
procedure Write_Port (S : Session; Value : Port_Value; Result : out Status);
procedure Write_Pin  (S : Session; Pin : Pin_Number; State : Pin_State;
                      Result : out Status);          --  read-modify-write
procedure Read_Port    (S : Session; Value : out Port_Value; Result : out Status);
procedure Read_Outputs (S : Session; Value : out Port_Value; Result : out Status);
procedure Set_Polarity (S : Session; Inverted : Port_Value; Result : out Status);
--  child: ESP32S3.TCA9555.Interrupts.Attach (Dev, Action) / Detach (Dev)
\end{lstlisting}

The example holds one \texttt{Session} for the whole test, forces every pin to an
input (so nothing fights the external wiring), then exercises the input port, the
output register round-trip, a per-pin RMW, and the polarity register.

\begin{lstlisting}
package GPX renames ESP32S3.TCA9555;
use type GPX.Status;  use type GPX.Port_Value;
Dev : GPX.Device;  S : GPX.Session;  St : GPX.Status;  V, Orig : GPX.Port_Value;
begin
   GPX.Setup (Dev, Addr => 0, Sda => 8, Scl => 7);
   GPX.Acquire (S, Dev);                           --  hold the expander

   GPX.Set_Directions (S, Inputs => 16#FFFF#, Result => St);  --  all inputs
   GPX.Read_Port (S, Orig, St);                    --  comms check

   GPX.Write_Port (S, 16#A55A#, St);               --  output register only
   GPX.Read_Outputs (S, V, St);                    --  V should read 0xA55A

   GPX.Write_Pin (S, 5, GPX.High, St);             --  per-pin RMW (Session-safe)
   GPX.Read_Outputs (S, V, St);

   GPX.Set_Polarity (S, 16#A55A#, St);
   GPX.Set_Polarity (S, 0, St);                    --  restore
   GPX.Release (S);
\end{lstlisting}

\noindent\textbf{Example:} \texttt{examples/esp32s3\_tca9555}.

\section{PCF85063A -- I2C Real-Time Clock}\index{PCF85063A driver}
\label{app:ext:pcf85063a}

The NXP PCF85063A is a small fixed-address (\texttt{Bus\_Address} = \texttt{0x51})
I2C real-time clock and calendar: seconds through years (stored BCD on the wire,
binary here), a programmable alarm matching any subset of
second/minute/hour/day/weekday, an oscillator-stop flag reporting whether time
survived a power loss, and an active-low open-drain \texttt{INT} line the alarm
asserts. Up to 400 kHz I2C.

\texttt{ESP32S3.PCF85063A} sits on \texttt{ESP32S3.I2C}. Unlike the expander
drivers it exposes no \texttt{Session} type to the caller -- each operation opens
its own short-lived controlled I2C session for one complete transaction and
auto-releases the host on scope exit, so concurrent callers serialise. Register
reads use the chip's auto-incrementing pointer (a 1-byte write sets it, a following
read streams from it). \texttt{Get\_Time} returns a \texttt{Valid} flag (clock
integrity); \texttt{Set\_Time} brackets its write with stop/start and clears the OS
flag. The optional \texttt{INT} pin is wired through \texttt{Setup} and armed by
\texttt{ESP32S3.PCF85063A.Interrupts.Attach} as a pulled-up falling-edge input. The
controlled session restricts this driver to the embedded and full profiles.

\begin{lstlisting}
Bus_Address : constant ESP32S3.I2C.Slave_Address := 16#51#;
subtype Year_Number is Natural range 2000 .. 2099;   --  also Month/Day/Hour/...
type Weekday is (Sunday, Monday, Tuesday, Wednesday, Thursday, Friday, Saturday);
type Time is record
   Year : Year_Number; Month : Month_Number; Day : Day_Number;
   Day_Of_Week : Weekday; Hour : Hour_Number;
   Minute : Minute_Number; Second : Second_Number;
end record;
type Alarm is record   --  any subset of fields participates
   Use_Second : Boolean; Second : Second_Number;  Use_Minute  : Boolean; ...
   Use_Hour, Use_Day, Use_Weekday : Boolean; ...
end record;
type Status is (OK, Bus_Error);   type Device is limited private;

procedure Setup (Dev : out Device; Sda, Scl : ESP32S3.GPIO.Pin_Id;
                 Int_Pin : ESP32S3.GPIO.Optional_Pin := ESP32S3.GPIO.No_Pin;
                 Host : ESP32S3.I2C.I2C_Host := ESP32S3.I2C.I2C0;
                 Clock_Hz : Positive := 400_000);
procedure Get_Time (Dev : Device; T : out Time; Valid : out Boolean;
                    Result : out Status);
procedure Set_Time (Dev : Device; T : Time; Result : out Status);
procedure Reset (Dev : Device; Result : out Status);
procedure Set_Alarm (Dev : Device; A : Alarm; Result : out Status);
procedure Alarm_Triggered (Dev : Device; Fired : out Boolean; Result : out Status);
procedure Acknowledge_Alarm (Dev : Device; Result : out Status);   --  releases INT
--  child: ESP32S3.PCF85063A.Interrupts.Attach (Dev, Action) / Detach (Dev)
\end{lstlisting}

The example probes, resets, loads a known calendar, arms a seconds-match alarm five
seconds out, then watches the clock tick. With no \texttt{INT} line wired
(\texttt{No\_Pin}) the alarm is found by polling \texttt{AF} over I2C;
\texttt{Attach} is then a no-op.

\begin{lstlisting}
package RTC renames ESP32S3.PCF85063A;
use type RTC.Status;
Dev : RTC.Device;  T : RTC.Time;  Valid : Boolean;  St : RTC.Status;
Initial : constant RTC.Time :=
  (Year => 2026, Month => 6, Day => 22, Day_Of_Week => RTC.Monday,
   Hour => 14, Minute => 30, Second => 0);
begin
   RTC.Setup (Dev, Sda => 8, Scl => 7, Int_Pin => ESP32S3.GPIO.No_Pin);
   RTC.Interrupts.Attach (Dev, Alarm_IRQ.Handler'Access);   --  no-op: No_Pin

   RTC.Get_Time (Dev, T, Valid, St);          --  probe (ACK?)
   RTC.Reset (Dev, St);
   RTC.Set_Time (Dev, Initial, St);           --  clears OS -> Valid = True
   RTC.Set_Alarm (Dev, (Use_Second => True, Second => 5, others => <>), St);

   for Tick in 1 .. 10 loop
      delay until Clock + Seconds (1);
      RTC.Get_Time (Dev, T, Valid, St);
      declare Fired : Boolean; begin
         RTC.Alarm_Triggered (Dev, Fired, St);
         if St = RTC.OK and then Fired then
            RTC.Acknowledge_Alarm (Dev, St);  --  releases INT
            exit;
         end if;
      end;
   end loop;
\end{lstlisting}

\noindent\textbf{Example:} \texttt{examples/esp32s3\_pcf85063a}.

\section{SHT41 -- I2C Temperature / Humidity Sensor}\index{SHT41 driver}
\label{app:ext:sht41}

The Sensirion SHT41 (SHT4x family) is a command-based I2C temperature and humidity
sensor with no registers: a measurement is ``write a 1-byte command, wait the
conversion time, read 6 bytes'' -- two CRC-8-protected 16-bit words (temperature
then humidity). The SHT41-AD1B answers at \texttt{Default\_Address} (\texttt{0x44});
other variants use \texttt{0x45}/\texttt{0x46}. There is no interrupt -- it is read
on request.

\texttt{ESP32S3.SHT41} sits on \texttt{ESP32S3.I2C} and shares the PCF85063A shape:
no caller-visible session, each operation opens a short-lived controlled I2C
session and auto-releases the host. \texttt{Measure} blocks for the conversion (per
\texttt{Precision}) while holding the bus and returns integer milli-units --
\texttt{Temperature} in milli-degrees Celsius, \texttt{Humidity} in milli-percent
RH -- so no float library is required; CRC failures surface as \texttt{CRC\_Error}.
The controlled session restricts this driver to the embedded and full profiles.

\begin{lstlisting}
Default_Address : constant ESP32S3.I2C.Slave_Address := 16#44#;
type Measurement is record
   Temperature : Integer := 0;   --  milli-degrees Celsius
   Humidity    : Integer := 0;   --  milli-percent RH (0 .. 100_000)
end record;
type Precision is (Low, Medium, High);          --  High ~8 ms, Low ~2 ms
type Status is (OK, Bus_Error, CRC_Error);
type Device is limited private;

procedure Setup (Dev : out Device; Sda, Scl : ESP32S3.GPIO.Pin_Id;
                 Address : ESP32S3.I2C.Slave_Address := Default_Address;
                 Host : ESP32S3.I2C.I2C_Host := ESP32S3.I2C.I2C0;
                 Clock_Hz : Positive := 400_000);
procedure Measure (Dev : Device; Value : out Measurement; Result : out Status;
                   Repeatability : Precision := High);
procedure Read_Serial_Number
  (Dev : Device; Serial : out Interfaces.Unsigned_32; Result : out Status);
procedure Reset (Dev : Device; Result : out Status);
\end{lstlisting}

The example reads the 32-bit serial number as a presence check, then takes a
high-precision measurement once a second.

\begin{lstlisting}
package SHT renames ESP32S3.SHT41;
use type SHT.Status;
Dev : SHT.Device;  St : SHT.Status;  SN : Interfaces.Unsigned_32;
begin
   SHT.Setup (Dev, Sda => 8, Scl => 7);

   SHT.Read_Serial_Number (Dev, SN, St);   --  doubles as a comms check
   Put ("serial : 0x");  Put_Hex (SN, 8);  New_Line;

   for Tick in 1 .. 15 loop
      delay until Clock + Seconds (1);
      declare M : SHT.Measurement; begin
         SHT.Measure (Dev, M, St);          --  blocks ~8 ms (High)
         exit when St /= SHT.OK;
         Put ("T=");  Put_Fixed (M.Temperature, 1000, 2);   --  m'C -> C
         Put (" C  RH=");  Put_Fixed (M.Humidity, 1000, 2);  Put_Line (" %");
      end;
   end loop;
\end{lstlisting}

\cnote{No vendor SDK is linked. The SHT4x's CRC-8 and the milli-unit fixed-point
conversion are done in Ada, so the whole sensor path is float-free -- unlike the
typical ESP-IDF driver that returns \texttt{float} degrees and pulls in soft-float.}

\noindent\textbf{Example:} \texttt{examples/esp32s3\_sht41}.

\section{QMI8658C -- I2C 6-Axis IMU}\index{QMI8658C driver}
\label{app:ext:qmi8658c}

The QST QMI8658C is a 6-axis IMU (3-axis accelerometer + 3-axis gyroscope) on I2C.
The \texttt{SA0}/\texttt{SDO} pin selects the 7-bit address:
\texttt{Address\_SA0\_High} (\texttt{0x6A}, the floating/power-on default via the
internal pull-up) or \texttt{Address\_SA0\_Low} (\texttt{0x6B}).
\texttt{WHO\_AM\_I} returns \texttt{Who\_Am\_I\_Value} (\texttt{0x05}). A
programmable \texttt{INT} line (INT1/INT2) is push-pull active-high by default.

\texttt{ESP32S3.QMI8658C} sits on \texttt{ESP32S3.I2C} with the PCF85063A shape (no
caller-visible session; short-lived auto-released I2C session per operation;
auto-incrementing register pointer set by \texttt{Configure}, run little-endian).
\texttt{Configure} sets the accelerometer and gyroscope full scale and a shared
output rate and enables both sensors; \texttt{Read\_Accelerometer}/%
\texttt{Read\_Gyroscope} return raw signed 16-bit counts in \texttt{Axes}, which
the \texttt{Accel\_LSB\_Per\_G}/\texttt{Gyro\_LSB\_Per\_DPS} accessors scale to
physical units. The optional \texttt{INT} pin is wired through \texttt{Setup} and
armed by \texttt{ESP32S3.QMI8658C.Interrupts.Attach} as a pulled-down rising-edge
input. The controlled session restricts this driver to the embedded and full
profiles.

\begin{lstlisting}
Address_SA0_High : constant ESP32S3.I2C.Slave_Address := 16#6A#;  --  default
Address_SA0_Low  : constant ESP32S3.I2C.Slave_Address := 16#6B#;
Who_Am_I_Value   : constant Interfaces.Unsigned_8     := 16#05#;
type Accel_Range is (Range_2G, Range_4G, Range_8G, Range_16G);
type Gyro_Range  is (Range_16DPS, ... Range_2048DPS);
type Output_Rate is (ODR_7520_Hz, ... ODR_235_Hz, ... ODR_29_Hz);
type Axes is record X, Y, Z : Interfaces.Integer_16 := 0; end record;
type Status is (OK, Bus_Error);   type Device is limited private;

procedure Setup (Dev : out Device; Sda, Scl : ESP32S3.GPIO.Pin_Id;
                 Int_Pin : ESP32S3.GPIO.Optional_Pin := ESP32S3.GPIO.No_Pin;
                 Address : ESP32S3.I2C.Slave_Address := Address_SA0_Low;
                 Host : ESP32S3.I2C.I2C_Host := ESP32S3.I2C.I2C0;
                 Clock_Hz : Positive := 400_000);
procedure Read_Who_Am_I (Dev : Device; Id : out Interfaces.Unsigned_8;
                         Result : out Status);
procedure Reset (Dev : Device; Result : out Status);     --  then wait ~15 ms
procedure Configure (Dev : in out Device; Accel : Accel_Range := Range_8G;
                     Gyro : Gyro_Range := Range_512DPS;
                     Rate : Output_Rate := ODR_235_Hz; Result : out Status);
procedure Read_Accelerometer (Dev : Device; A : out Axes; Result : out Status);
procedure Read_Gyroscope     (Dev : Device; G : out Axes; Result : out Status);
procedure Read_Temperature
  (Dev : Device; Raw : out Interfaces.Integer_16; Result : out Status);
function  Accel_LSB_Per_G  (Dev : Device) return Positive;
function  Gyro_LSB_Per_DPS (Dev : Device) return Positive;
\end{lstlisting}

The example probes both SA0 addresses for \texttt{WHO\_AM\_I} = \texttt{0x05},
resets and configures the part, then streams accel and gyro in milli-units using
the sensitivity accessors and integer math.

\begin{lstlisting}
package IMU renames ESP32S3.QMI8658C;
use type IMU.Status;
Dev : IMU.Device;  St : IMU.Status;  Id : Unsigned_8;  Found : Boolean := False;
Candidates : constant array (1 .. 2) of ESP32S3.I2C.Slave_Address :=
  (IMU.Address_SA0_Low, IMU.Address_SA0_High);
begin
   for I in Candidates'Range loop
      IMU.Setup (Dev, Sda => 8, Scl => 7, Address => Candidates (I));
      IMU.Read_Who_Am_I (Dev, Id, St);
      if St = IMU.OK and then Id = IMU.Who_Am_I_Value then
         Found := True;  exit;
      end if;
   end loop;

   IMU.Reset (Dev, St);
   delay until Clock + Milliseconds (15);          --  reset settle
   IMU.Configure (Dev, Accel => IMU.Range_8G, Gyro => IMU.Range_512DPS,
                  Rate => IMU.ODR_235_Hz, Result => St);

   declare
      A, G  : IMU.Axes;
      A_Lsb : constant Positive := IMU.Accel_LSB_Per_G (Dev);
   begin
      IMU.Read_Accelerometer (Dev, A, St);
      IMU.Read_Gyroscope (Dev, G, St);
      --  milli-g = Integer (A.X) * 1000 / A_Lsb
   end;
\end{lstlisting}

\noindent\textbf{Example:} \texttt{examples/esp32s3\_qmi8658c}.

\section{ES8311 -- Audio Codec}\index{ES8311 driver}
\label{app:ext:es8311}

The ES8311 is an Everest low-power \emph{mono} audio codec: a register-configured
DAC/ADC pair that the SoC controls over I2C while carrying PCM audio over I2S. The
driver in \texttt{ESP32S3.ES8311} brings it up for 16-bit output, running the
codec-vendor register-init sequence on I2C and configuring an I2S port as the
master that generates MCLK, BCLK (SCLK) and LRCK and shifts PCM out on the codec's
DSDIN line. The codec runs as an I2S slave clocked from the ESP's MCLK $= 256
\times$ sample-rate (4.096\,MHz at 16\,kHz); the register coefficients are
rate-independent for that $256\times$/16-bit configuration.

It sits on two on-chip drivers at once -- \texttt{ESP32S3.I2C} for control and
\texttt{ESP32S3.I2S} for data. \texttt{Setup} is a one-time, single-threaded
bring-up; pass an \texttt{Asdout} pin to also enable the ADC (microphone) capture
path. After setup, audio is owned through a limited, controlled \texttt{Output}
handle that holds the I2S port exclusively and releases it on scope exit, so it
requires a tasking runtime (embedded/full profile). \texttt{Play\_Continuous}
starts a self-looping DMA for gapless, click-free, zero-CPU playback;
\texttt{Capture} samples the ADC underneath it without disturbing playback or the
shared master clock.

\begin{lstlisting}
package ESP32S3.ES8311 is
   subtype Address is ESP32S3.I2C.Slave_Address;
   Default_Address : constant Address := 16#18#;

   procedure Setup
     (I2C_Bus     : ESP32S3.I2C.I2C_Host;
      Sda, Scl    : ESP32S3.GPIO.Pin_Id;
      Port        : ESP32S3.I2S.I2S_Port;
      Mclk, Sclk, Lrck, Dsdin : ESP32S3.GPIO.Pin_Id;
      Asdout      : ESP32S3.GPIO.Optional_Pin := ESP32S3.GPIO.No_Pin;
      Sample_Rate : Positive := 16_000;
      Volume      : Natural  := 70;     --  DAC volume, 0 .. 100 %
      Mic_Gain_Db : Natural  := 24;     --  ADC PGA gain, 0 .. 42 dB
      Addr        : Address  := Default_Address;
      Ok          : out Boolean);

   procedure Set_Volume (Percent : Natural; Ok : out Boolean);

   type Output is limited private;          --  exclusive, auto-releasing
   Not_Ready : exception;

   procedure Acquire (O : in out Output);
   procedure Play            (O : Output; Samples : System.Address; Length : Natural);
   procedure Play_Continuous (O : Output; Samples : System.Address; Length : Natural);
   procedure Capture         (O : Output; Samples : System.Address; Length : Natural);
   procedure Stop    (O : Output);
   procedure Release (O : in out Output);
end ESP32S3.ES8311;
\end{lstlisting}

The example plays a gapless 440\,Hz sine on the DAC and simultaneously captures the
ADC to estimate the tone heard back -- a loopback test. Control is on I2C (SDA=IO8,
SCL=IO7); audio on I2S (MCLK=IO1, SCLK=IO2, LRCK=IO4, DSDIN=IO5 out, ASDOUT=IO3 in).

\begin{lstlisting}
declare
   Ok        : Boolean;
   Audio     : ESP32S3.ES8311.Output;
begin
   ESP32S3.ES8311.Setup
     (I2C_Bus => ESP32S3.I2C.I2C0,
      Sda     => 8, Scl => 7,
      Port    => ESP32S3.I2S.I2S0,
      Mclk    => 1, Sclk => 2, Lrck => 4, Dsdin => 5,
      Asdout  => 3,                       --  codec ADC out -> our data in
      Sample_Rate => 16_000, Volume => 75, Mic_Gain_Db => 24, Ok => Ok);

   ESP32S3.ES8311.Acquire (Audio);        --  hold the audio port
   --  Self-looping DMA: replays Buf forever, click-free, zero CPU.
   ESP32S3.ES8311.Play_Continuous (Audio, Buf'Address, Buf_Bytes);
   loop                                   --  capture under the running tone
      ESP32S3.ES8311.Capture (Audio, Cap'Address, Cap_Bytes);
      Analyse (Peak, Est);                --  should read ~440 Hz
      delay until Clock + Milliseconds (1000);
   end loop;
end;
\end{lstlisting}

\noindent\textbf{Example:} \texttt{examples/esp32s3\_es8311}.

\section{ST7789 -- SPI TFT Display Controller}\index{ST7789 driver}
\label{app:ext:st7789}

The ST7789 (ST77xx family) drives a small colour TFT panel over 4-wire SPI. The
link is \emph{write-only}: CLK + MOSI, no MISO, so the controller cannot be read
back and there is no probe or status query. Three GPIOs are driven directly -- DC
(data/command), CS (chip select), and an optional RST (software reset is used if it
is \texttt{No\_Pin}). Pixels are 16-bit RGB565, MSB-first. \texttt{ESP32S3.ST7789}
sits on \texttt{ESP32S3.SPI} (mode 0, full-duplex master, default 40\,MHz); CS is
driven as a plain GPIO rather than routed to the peripheral.

Locking is two-level, like the on-chip expander drivers. A \texttt{Session} is an
exclusive, RAII hold on one display, taken with \texttt{Acquire} and auto-released
on scope exit; hold it across a whole sequence so no other task corrupts the
controller mid-stream. The SPI host itself is locked only inside each operation
(assert CS, transfer, deassert CS, release), so the bus is free between operations.
The per-display guards are a library-level array keyed by the CS pin, and the
controlled \texttt{Session} means embedded/full profiles only.

A companion child \texttt{ESP32S3.ST7789.Text} layers a 5$\times$7 bitmap font on
the parent's \texttt{Draw\_Bitmap} primitive: \texttt{Draw\_Char}/\texttt{Draw\_Text}
blit opaque 6$\times$8 cells (the panel is write-only, so there is no read-back for a
transparent overlay) at an integer \texttt{Scale}, with LF handling for newlines.

\begin{lstlisting}
package ESP32S3.ST7789 is
   type Color is mod 2 ** 16;                        --  RGB565, MSB-first
   function RGB (R, G, B : Natural) return Color;    --  each 0 .. 255
   Black, White, Red, Green, Blue : constant Color;  --  named colours
   type Color_Array is array (Natural range <>) of Color;
   type Rotation is (Rot_0, Rot_90, Rot_180, Rot_270);

   type Device  is limited private;
   type Session is limited private;
   Not_Initialized, Not_Owned : exception;

   procedure Setup
     (Dev                : out Device;
      Sclk, Mosi, DC, CS : ESP32S3.GPIO.Pin_Id;
      Width, Height      : Positive := 240;
      RST                : ESP32S3.GPIO.Optional_Pin := ESP32S3.GPIO.No_Pin;
      Clock_Hz           : Positive := 40_000_000);

   procedure Acquire (S : in out Session; Dev : Device);
   procedure Release (S : in out Session);

   procedure Init         (S : Session);             --  reset + power-on sequence
   procedure Set_Rotation (S : Session; Rot : Rotation);
   procedure Fill         (S : Session; C : Color);
   procedure Fill_Rect    (S : Session; X, Y, W, H : Natural; C : Color);
   procedure Set_Pixel    (S : Session; X, Y : Natural; C : Color);
   procedure Draw_Bitmap  (S : Session; X, Y, W, H : Natural; Pixels : Color_Array);
end ESP32S3.ST7789;
\end{lstlisting}

The example drives a 240$\times$240 panel on SPI2 (SCLK=IO12, MOSI=IO13, DC=IO16,
CS=IO10; backlight on IO6 is driven by the example, not the driver). One
\texttt{Session} is held for the whole demo while each fill locks the bus only for
its own transfer.

\begin{lstlisting}
declare
   Dev : ESP32S3.ST7789.Device;
   S   : ESP32S3.ST7789.Session;
begin
   ESP32S3.ST7789.Setup (Dev, Sclk => 12, Mosi => 13, DC => 16, CS => 10);
   ESP32S3.ST7789.Acquire (S, Dev);            --  protect the display
   ESP32S3.ST7789.Init (S);
   ESP32S3.ST7789.Fill (S, ESP32S3.ST7789.Red);
   ESP32S3.ST7789.Fill_Rect (S, X => 70, Y => 70, W => 100, H => 100,
                             C => ESP32S3.ST7789.White);
   ESP32S3.ST7789.Text.Draw_Text
     (S, X => 6, Y => 16, Str => "ESP32-S3 + Ada",
      FG => ESP32S3.ST7789.White, BG => ESP32S3.ST7789.RGB (0, 0, 32));
   ESP32S3.ST7789.Release (S);
end;                                            --  Session auto-released
\end{lstlisting}

\noindent\textbf{Example:} \texttt{examples/esp32s3\_st7789}.

\section{TLV2556 -- SPI 12-bit ADC}\index{TLV2556 driver}\index{ADC driver}
\label{app:ext:tlv2556}

The TLV2556 is a Texas Instruments 12-bit, 11-channel, 200-kSPS
successive-approximation ADC with an internal reference, on an SPI-compatible
serial interface (mode 0). It is the off-chip complement to the SoC's own SAR ADC
(in the HAL chapter): more channels, a stable internal reference, and -- because
it lives on the SPI bus -- usable even when the internal ADC's pins are committed
elsewhere.

Two things make it a useful worked example. First, it \emph{shares the SPI2 bus
with the W25Q flash}, through the SPI driver's built-in software chip-select: you
name the select GPIO as \texttt{CS\_Pin} and the driver drives it -- active-low,
held across the whole conversion -- and suppresses the hardware CS0 for the hold,
so it cannot disturb the flash. (Contrast the ST7789 above, which toggles its CS
as a plain GPIO by hand.) Second, the converter is \emph{pipelined}: each 16-clock
frame shifts an 8-bit command in (channel select plus configuration) while the
\emph{previous} conversion's result shifts out, MSB-first. \texttt{Read} hides
that -- it primes the conversion, waits out the (max $\sim$5.5\,\textmu s)
conversion time, then reads the result back.

\begin{lstlisting}
package ESP32S3.TLV2556 is
   type Sample       is range 0 .. 4095;              --  12-bit unipolar result
   type Analog_Input is (AIN0, AIN1, AIN2, AIN3, AIN4, AIN5, AIN6, AIN7,
                         AIN8, AIN9, AIN10,            --  the 11 analog inputs ...
                         Test_Half, Test_Zero, Test_Full);  --  ... and 3 self-tests
   type Reference    is (Internal_4096mV, Internal_2048mV, External);

   type Device is record                              --  describes itself on the bus
      Host     : ESP32S3.SPI.SPI_Host;
      Clock_Hz : Positive                  := 8_000_000;
      CS_Pin   : ESP32S3.GPIO.Optional_Pin := ESP32S3.GPIO.No_Pin;  -- driver-driven CS
      CS_CB    : ESP32S3.SPI.CS_Select     := null;   --  or a callback (decoder/expander)
      Ctx      : System.Address            := System.Null_Address;
   end record;

   procedure Initialize (Dev : Device; Ref : Reference := External);
   function  Read       (Dev : Device; Input : Analog_Input) return Sample;
   function  Millivolts (S : Sample; Ref : Reference) return Natural;
end ESP32S3.TLV2556;
\end{lstlisting}

The standout feature for bring-up is the \emph{self-test inputs}: three of the
\texttt{Analog\_Input} values are not pins at all but internal voltages
\emph{ratiometric to the reference rails}, so \texttt{Test\_Zero} reads 0,
\texttt{Test\_Half} reads 2048, and \texttt{Test\_Full} reads 4095 -- regardless
of the reference voltage or any analog wiring. They are the ADC analog of the
W25Q's JEDEC-ID check: a reference- and wiring-independent proof that the SPI link,
the framing, and the converter all work, before you trust a real reading.

\begin{lstlisting}
declare
   package ADC renames ESP32S3.TLV2556;
   Dev : ADC.Device :=                                --  shares SPI2; CS on IO12
     (Host => ESP32S3.SPI.SPI2, Clock_Hz => 8_000_000, CS_Pin => 12, others => <>);
begin
   ESP32S3.SPI.Setup (ESP32S3.SPI.SPI2);              --  bus shared with the flash
   ESP32S3.SPI.Configure_Pins (ESP32S3.SPI.SPI2, Sclk => 1, Mosi => 4, Miso => 45);
   ADC.Initialize (Dev, Ref => ADC.External);

   --  Bring-up proof: the ratiometric self-tests, independent of reference/wiring.
   Put_Line ("zero=" & ADC.Read (Dev, ADC.Test_Zero)'Image    -- ~ 0
           & " half=" & ADC.Read (Dev, ADC.Test_Half)'Image   -- ~ 2048
           & " full=" & ADC.Read (Dev, ADC.Test_Full)'Image); -- ~ 4095

   Value := ADC.Read (Dev, ADC.AIN0);                 --  then a real channel
end;
\end{lstlisting}

On the hardware (the self-tests carry a couple of LSB of analog slack):

\begin{lstlisting}[style=sh]
[tlv2556] self-test: zero=0 half=2047 full=4095   PASS
[tlv2556] AIN0 = 195 / 4095
\end{lstlisting}

\noindent\textbf{Example:} \texttt{examples/esp32s3\_tlv2556}.

\section{GPS -- NMEA-0183 Receiver}\index{GPS driver}
\label{app:ext:gps}

The GPS driver targets a UART GNSS module (developed against the Quectel L76K) that
streams NMEA-0183 sentences. Unlike the passive I2C device drivers you poll through
a handle, \texttt{ESP32S3.GPS} is a singleton background \emph{service}: a
library-level task owns one UART, continuously reads and decodes the sentence
stream, and publishes results into a protected store. The application has no
\texttt{Device} handle -- it just reads the store. It sits on \texttt{ESP32S3.UART};
only \texttt{Rx} is required (wired to the module's TXD), with optional \texttt{Tx}
(for configuration) and \texttt{Pps} (a 1PPS GPIO interrupt for time alignment).

Every published value is written and read under the store's lock, so a reader never
sees a half-updated value; in particular \texttt{Latitude}/\texttt{Longitude} are
one \texttt{Position} record updated atomically. Each reading carries the
\texttt{Updated\_At} time at which a \emph{valid} sentence last refreshed it --
compare \texttt{Age} against your tolerance to detect staleness. The driver decodes
GGA, RMC, GLL, VTG, GSV, GSA and ZDA into fix, position, time, date, velocity,
signal and satellite-in-view readings. The task plus protected objects mean
embedded/full profiles only.

\cnote{Latitude and longitude follow the u-blox convention: signed integers in
$10^{-7}$ degrees (\texttt{Integer\_32}), not floats -- exact, atomically paired,
and free of any soft-float dependency in the bare runtime.}

A child \texttt{ESP32S3.GPS.L76K} adds the Quectel PCAS vendor commands
(constellation selection, baud, update rate, restart) sent verbatim through
\texttt{Send}; only \texttt{Set\_Constellation} (PCAS04) is hardware-tested.
\texttt{Inject} feeds one framed sentence straight into the decoder for
host/on-target self-test with no live receiver.

\begin{lstlisting}
package ESP32S3.GPS is
   type Position is record                       --  1e-7 degrees, +N/-S +E/-W
      Latitude  : Interfaces.Integer_32 := 0;
      Longitude : Interfaces.Integer_32 := 0;
   end record;
   type Fix_Quality is (No_Fix, GPS_Fix, DGPS_Fix);
   type Fix_Type    is (Fix_None, Fix_2D, Fix_3D);

   type Position_Reading is record               --  value + freshness + valid
      Value : Position; Updated_At : Ada.Real_Time.Time; Valid : Boolean := False;
   end record;
   --  Fix_Reading, Time_Reading, Date_Reading, Velocity_Reading, Signal_Reading ...

   procedure Setup
     (Port : ESP32S3.UART.UART_Port;
      Rx   : ESP32S3.GPIO.Optional_Pin;
      Tx   : ESP32S3.GPIO.Optional_Pin := ESP32S3.GPIO.No_Pin;
      Pps  : ESP32S3.GPIO.Optional_Pin := ESP32S3.GPIO.No_Pin;
      Baud : ESP32S3.UART.Baud_Rate    := 9600);

   function Current_Position return Position_Reading;
   function Current_Fix      return Fix_Reading;
   function Current_Time     return Time_Reading;
   function Age (Updated_At : Ada.Real_Time.Time) return Ada.Real_Time.Time_Span;
   procedure Inject (Sentence : String; Accepted : out Boolean);
end ESP32S3.GPS;
\end{lstlisting}

The example self-tests the decoder by injecting canned sentences (including a
bad-checksum reject) before \texttt{Setup}, then brings up UART0 (Rx=IO44, Tx=IO43,
9600 baud) and prints the latest fix and its age once a second.

\begin{lstlisting}
GPS.Inject ("$GPGGA,123519,4807.038,N,01131.000,E,1,08,0.9,545.4,M,46.9,M,,*47", Ok);
declare
   P : constant GPS.Position_Reading := GPS.Current_Position;
begin
   pragma Assert (P.Valid and then P.Value.Latitude = 481_173_000);
end;

GPS.Setup (Port => ESP32S3.UART.UART0, Rx => 44, Tx => 43, Baud => 9_600);
loop
   delay until Clock + Seconds (1);
   declare
      P : constant GPS.Position_Reading := GPS.Current_Position;
      Fresh : constant Boolean :=
        P.Valid and then To_Duration (GPS.Age (P.Updated_At)) < 3.0;
   begin
      null;   --  print P.Value when Fresh, else show acquisition progress
   end;
end loop;
\end{lstlisting}

\noindent\textbf{Example:} \texttt{examples/esp32s3\_gps}.

\section{W25Q -- SPI NOR Flash}\index{W25Q driver}\index{SPI flash}
\label{app:ext:w25q}

The W25Q256FV is a 32\,MB (256\,Mbit) SPI NOR flash -- the on-board non-volatile
store the local filesystem of Chapter~\ref{ch:storage} lives on. It sits on SPI2,
sharing the bus with the W5500 (and, in the ADC example, the TLV2556) through the
SPI driver's software chip-select. The driver is deliberately small: identify, read,
erase by sector, program by page -- the primitives a wear-levelling layer and a
filesystem build on.

\begin{lstlisting}
package ESP32S3.W25Q is
   Page_Size   : constant := 256;        --  program granularity (one Page-Program)
   Sector_Size : constant := 4096;       --  smallest erase unit

   type JEDEC_ID is record
      Manufacturer : Unsigned_8;         --  0xEF for Winbond
      Memory_Type  : Unsigned_8;
      Capacity     : Unsigned_8;         --  log2 of the size in bytes
   end record;

   type Flash is record                  --  describes itself on the bus
      Host     : ESP32S3.SPI.SPI_Host;
      Clock_Hz : Positive                  := 8_000_000;
      CS_Pin   : ESP32S3.GPIO.Optional_Pin := ESP32S3.GPIO.No_Pin;  -- driver-driven CS
      CS_CB    : ESP32S3.SPI.CS_Select     := null;   --  or a callback
      Ctx      : System.Address            := System.Null_Address;
   end record;

   procedure Read_Identification (Dev : Flash; ID : out JEDEC_ID);
   function  Capacity_Bytes      (ID : JEDEC_ID) return Address;
   procedure Initialize          (Dev : Flash; OK : out Boolean);

   procedure Read         (Dev : Flash; Addr : Address; Data : out Byte_Array);
   procedure Erase_Sector (Dev : Flash; Addr : Address);
   procedure Program_Page (Dev : Flash; Addr : Address; Data : Byte_Array)
     with Pre => Data'Length in 1 .. Page_Size
                 and then Natural (Addr mod Page_Size) + Data'Length <= Page_Size;
end ESP32S3.W25Q;
\end{lstlisting}

The bring-up proof is the \emph{JEDEC-ID check}: \texttt{Read\_Identification}
returns \texttt{EF 40 19} on a healthy W25Q256 (Winbond, 256\,Mbit) -- a
wiring-independent confirmation that the SPI link and framing work before any read
or write is trusted, the flash analogue of the TLV2556's self-test inputs
(\S\ref{app:ext:tlv2556}). The page/sector asymmetry is the usual NOR contract: you
may program any aligned 256-byte page but can only erase a whole 4\,KB sector, and
the \texttt{Program\_Page} precondition enforces the page rule at the call site.
Above this driver, \texttt{Block\_Dev.WL} adds wear levelling and \texttt{Ext4.FS} a
filesystem (Chapter~\ref{ch:storage}); from the LISP REPL, \texttt{(spi-xfer ...)}
reads the same JEDEC id by hand (Chapter~\ref{ch:lisp}).

\noindent\textbf{Example:} \texttt{examples/esp32s3\_w25q}.

\section{SD Card (SD\_SPI) -- Removable Storage}\index{SD card driver}
\label{app:ext:sdspi}

\texttt{ESP32S3.SD\_SPI} talks to an SD/SDHC card in SPI mode, layering the SD
command protocol (CMD0/8/58, ACMD41, CMD17/24, CRC7) on top of the task-safe
\texttt{ESP32S3.SPI} master. Four lines plus power are routed (MOSI/MISO/SCLK and a
CS driven as a plain GPIO, so it can stay asserted across a whole
command/response/data sequence); a 10\,k pull-up on MISO is recommended. Cards
initialise at $\leq$400\,kHz then run at \texttt{Data\_Clock\_Hz}. The API is always
512-byte LBA blocks; SDHC/SDXC block-addressing vs.\ older byte-addressing is
handled internally. Every operation takes the SPI Session for the whole transaction
(so concurrent callers serialise), which means embedded/full profiles only.

\emph{(Covered in full in the Mass Storage chapter, \S\ref{ch:storage}.)} This
entry is just the bus-level driver beneath that chapter's filesystem stack.

\begin{lstlisting}
package ESP32S3.SD_SPI is
   type Block         is array (0 .. 511) of Interfaces.Unsigned_8;
   type Block_Address is new Interfaces.Unsigned_32;
   type Card_Kind is (Unknown, SD_V1, SD_V2_SC, SD_V2_HC);
   type Status    is (OK, No_Card, Unusable, Init_Timeout, Read_Error, Write_Error);
   type Card is limited private;

   procedure Setup (C : out Card; Host : ESP32S3.SPI.SPI_Host;
                    Sclk, Mosi, Miso, Cs : ESP32S3.GPIO.Pin_Id;
                    Init_Clock_Hz : Positive := 400_000;
                    Data_Clock_Hz : Positive := 8_000_000);
   procedure Initialize  (C : in out Card; Result : out Status);
   function  Kind        (C : Card) return Card_Kind;
   procedure Read_Block  (C : in out Card; LBA : Block_Address;
                          Data : out Block; Result : out Status);
   procedure Write_Block (C : in out Card; LBA : Block_Address;
                          Data : Block; Result : out Status);
end ESP32S3.SD_SPI;
\end{lstlisting}

The example does a non-destructive round-trip on a scratch sector on SPI2
(SCLK=IO12, MOSI=IO11, MISO=IO13, CS=IO10): read it, write the same bytes back,
re-read, and verify they match.

\begin{lstlisting}
ESP32S3.SD_SPI.Setup (C, Host => ESP32S3.SPI.SPI2,
                      Sclk => 12, Mosi => 11, Miso => 13, Cs => 10);
ESP32S3.SD_SPI.Initialize (C, St);            --  Kind (C) now reports the card type
if St = ESP32S3.SD_SPI.OK then
   ESP32S3.SD_SPI.Read_Block  (C, Test_LBA, Orig, St);
   ESP32S3.SD_SPI.Write_Block (C, Test_LBA, Orig, St);   --  same bytes back
   ESP32S3.SD_SPI.Read_Block  (C, Test_LBA, Back, St);
   pragma Assert (St = ESP32S3.SD_SPI.OK and then Orig = Back);
end if;
\end{lstlisting}

\noindent\textbf{Example:} \texttt{examples/esp32s3\_sd\_spi}.

\section{TX1812 -- Addressable RGB LED String}\index{TX1812 driver}
\label{app:ext:tx1812}

The TX1812 is a single-wire, daisy-chainable RGB LED in the WS2812 (``NeoPixel'')
family: 24 bits of colour are clocked in MSB-first as a precisely timed pulse train
(a `1' is long-high/short-low, a `0' short-high/long-low), and a $>$80\,\textmu s
low ``reset'' latches the chain. \texttt{ESP32S3.TX1812} generates that waveform
with the \texttt{ESP32S3.RMT} peripheral, one RMT symbol per data bit. A
\texttt{Strip} is a claimed handle that takes an RMT transmit channel on
\texttt{Acquire} and releases it on scope exit (its RMT channel component is
controlled).

A \texttt{Strip} is statically sized by its \texttt{Count} discriminant: it carries
both the \texttt{Count}-pixel colour buffer and the \texttt{Count}$\times$24
RMT-symbol frame buffer, so declaring e.g.\ \texttt{Panel : Strip (64)} reserves all
of it at elaboration with no heap, and the linker verifies it fits. Set colours into
the buffer with \texttt{Set}/\texttt{Set\_All}, then \texttt{Show} clocks them out;
longer strings are streamed via RMT wrap re-fill, so the default of one RMT block
handles any length.

\begin{lstlisting}
package ESP32S3.TX1812 is
   type Color is record R, G, B : Unsigned_8 := 0; end record;
   Off, Red, Green, Blue, White : constant Color;

   --  Statically sized: holds the pixel buffer AND the Count*24 RMT frame.
   type Strip (Count : Positive) is limited private;

   procedure Acquire (S       : in out Strip;
                      Pin     : ESP32S3.GPIO.Pin_Id;
                      Channel : ESP32S3.RMT.TX_Index := 0;
                      Blocks  : Positive := 1);
   function  Is_Valid (S : Strip) return Boolean;
   procedure Release  (S : in out Strip);
   procedure Set      (S : in out Strip; Index : Positive; C : Color);
   procedure Set_All  (S : in out Strip; C : Color);
   procedure Show     (S : in out Strip);   --  clock the buffer out and latch
end ESP32S3.TX1812;
\end{lstlisting}

The example drives a 64-LED string on IO48 (RMT channel 0), declared once at
library level so its $\sim$6.4\,KiB of buffers are reserved at elaboration. With no
physical string attached, the board's on-board LED on IO48 is pixel 1 of the chain,
so it cycles red\,$\rightarrow$\,green\,$\rightarrow$\,blue\,$\rightarrow$\,white\,$\rightarrow$\,off.

\begin{lstlisting}
--  LED_Panel.Panel : ESP32S3.TX1812.Strip (64);   (library-level)
LED.Acquire (LED_Panel.Panel, Pin => 48, Channel => 0);
if not LED.Is_Valid (LED_Panel.Panel) then
   loop delay until Clock + Seconds (3600); end loop;   --  channel busy
end if;
loop
   for I in Colors'Range loop
      LED.Set_All (LED_Panel.Panel, Colors (I));   --  all 64 pixels
      LED.Show    (LED_Panel.Panel);               --  stream the whole frame
      delay until Clock + Milliseconds (600);
   end loop;
end loop;
\end{lstlisting}

\noindent\textbf{Example:} \texttt{examples/esp32s3\_tx1812}.

\section{W5500 -- SPI Ethernet Controller}\index{W5500 driver}
\label{app:ext:w5500}

The WIZnet W5500 is an Ethernet controller with a hardwired TCP/IP stack (ARP,
ICMP, IP, TCP, UDP in silicon), eight hardware sockets and 32\,KB of buffer RAM,
driven over SPI. On this board it sits on \texttt{ESP32S3.SPI} (SPI2) with
\texttt{MISO=IO45}, \texttt{MOSI=IO4}, \texttt{SCLK=IO1}, \texttt{SCSn=IO39},
\texttt{INTn=IO3} and \texttt{RSTn=IO11}; the chip-select is an ordinary GPIO held
low across each variable-length SPI frame.

Unlike the other external parts, the W5500 is a whole subsystem rather than a single
driver: a transport layer (\texttt{ESP32S3.W5500}), a socket engine
(\texttt{ESP32S3.W5500.Sockets}), optional interrupts, a DHCP client, and an adapter
(\texttt{ESP32S3.W5500.Net\_Device}) that presents the chip as a standard
\texttt{GNAT.Sockets} interface. Because that is far more than an API digest, it has
its own chapter -- see chapter~\ref{ch:networking}, ``Networking: GNAT.Sockets over
the W5500,'' for the driver layers, the chip-neutral \texttt{Net\_Devices}
interface, the DNS and DHCP helpers, and the worked examples (echo server, HTTP,
NTP, DNS, DHCP, and a weather client). Like the other controlled-handle drivers it
requires the embedded or full profile.

\noindent\textbf{Examples:} \texttt{examples/esp32s3\_w5500*}.
